%\documentclass[11pt]{amsart}
\documentclass[12pt,letterpaper]{report}
% Font and Encoding
\usepackage[T1]{fontenc}
\usepackage{times}  % Times Roman font
% Geometry for page layout
\usepackage[letterpaper,top=1in, bottom=1.25in, left=1.25in, right=1in, headheight=0.5in, headsep=0.25in]{geometry}
% Graphics and Images
\usepackage{graphicx}
\usepackage{epstopdf}
\DeclareGraphicsRule{.tif}{png}{.png}{`convert #1 `dirname #1`/`basename #1 .tif`.png}
% Math Packages
\usepackage{amsmath}
\usepackage{amssymb}
\usepackage{amsfonts}  % For additional math symbols
\usepackage{bm}        % Bold math symbols
% Text Formatting and Layout
\usepackage{xcolor}
\usepackage{setspace}    % Allows setting line spacing
\usepackage{float}       % Improved placement of figures
\usepackage{adjustbox}   % Optional for extra control on graphics
\usepackage{makecell}    % Advanced table cell formatting
% Tables and Columns
\usepackage{tabularx,booktabs}
\usepackage{multirow}
\usepackage{multicol}
\usepackage{ulem}
% Listings for Code
\usepackage{listings}
% Algorithms
\usepackage{algorithm}
\usepackage{algorithmic}
% URLs and Hyperlinks
\usepackage{url}         % URL formatting
\usepackage[hidelinks]{hyperref}  % Clickable links in PDF
% Bibliography
\usepackage{natbib}      % BibTeX bibliography style
% Captions and TOC
\usepackage[labelfont=bf]{caption}
\usepackage[titles]{tocloft}
% Additional Packages
\usepackage[section]{placeins}  % Place figures/tables within sections
\usepackage{titlesec}        % Customizable section titles
\usepackage{enumitem}        % Enhanced list formatting
\usepackage{afterpage}       % Code to run after a page break
\usepackage{pdflscape}       % Landscape pages
\usepackage{fancyhdr}        % Fancy headers and footers
%\usepackage{amsaddr}         % AMS address formatting
\usepackage{authblk}
\usepackage{lineno}          % Line numbers
\usepackage{newfloat}        % Define new floating environments
%\usepackage{mdsymbol}      % Conflicts with mathptmx, comment out if needed
% TikZ for Diagrams
\usepackage{tikz}
% CSV Simple for Reading CSV Files
\usepackage{csvsimple}


\newcommand{\fixedfigwidth}{0.95\textwidth}  % Define a consistent figure width

\newcolumntype{P}[1]{>{\centering\arraybackslash}p{#1}}
\newcolumntype{M}[1]{>{\centering\arraybackslash}m{#1}}

\newcommand{\betaoneone}{\beta_{11}}
\newcommand{\betaonetwo}{\beta_{12}}
\newcommand{\betatwoone}{\beta_{21}}
\newcommand{\betatwotwo}{\beta_{22}}

\newcommand{\boldGamma}{\boldsymbol \Gamma}
\newcommand{\boldU}{\boldsymbol U}


% Define a new command \address{}
\newcommand{\address}[1]{%
	\begin{center}
		\large\textit{#1}
	\end{center}
}


\DeclareFloatingEnvironment[name={Supplementary Figure},fileext=lsf,listname={List of Supplementary Figures}]{suppfigure}
\DeclareFloatingEnvironment[name={Supplementary Table},fileext=lsf,listname={List of Supplementary Tables}]{supptable}
\DeclareFloatingEnvironment[name={Supplementary Material},fileext=lsf,listname={List of Supplementary Material}]{suppmat}
\DeclareFloatingEnvironment[name={Supplementary Text},fileext=lsf,listname={List of Supplementary Text}]{supptext}

\linenumbers

\DeclareMathOperator*{\argmax}{arg\,max}
\DeclareMathOperator*{\argmin}{arg\,min}
\newcommand{\indep}{\perp \!\!\! \perp}
\newcommand{\nindep}{\not\!\perp\!\!\!\perp}

\newcommand{\boldbeta}{\boldsymbol \beta}

\renewcommand\thesection{\arabic{section}}

\begin{document}
	\title{Inferring causal networks for genomic trios with high-dimensional confounding variables}
	\author[1,2]{Jarred M. Kvamme}
	\author [3,4,5$\dagger$]{Audrey Qiuyan Fu}
	
	\affil[1]{Research Computing and Data Services, University of Idaho}
	\affil[2]{Institute for Interdisciplinary Data Sciences, University of Idaho}
	\affil[3]{Department of Family Medicine and Public Health Sciences, Wayne State University School of Medicine}
	\affil[4]{Center for Molecular Medicine and Genetics, Wayne State University School of Medicine}
	\affil[5]{Institute for AI and Data Science, Wayne State University}
	\affil[$\dagger$]{Correspondence: \href{audreyfu@wayne.edu}{\textit{audreyfu@wayne.edu}}}
	\maketitle
	
\begin{doublespace}

%\begin{abstract}
\section*{Abstract}
%Network inference under the principle of Mendelian randomization makes it possible to infer the potential causal relationship between two biological variables (e.g., expression of two genes), using a genetic variant as an instrumental variable and thus forming a trio.  Inference is often influenced by many other variables (such as other genes and demographic variables).  Being able to properly account for confounding is therefore essential to causal network inference.  However,  methods that can deal with many confounding variables are limited.  Here, we introduce a regression-based method under this principle, namely Mendelian Randomization for Genomic Networks (MRGN), to detect diverse causal networks for genomic trios, beyond the canonical mediation model, using individual-level data while accounting for many confounding variables.  
%We demonstrate the inference power, robustness and computational efficiency through extensive simulation and application to the Genotype-Tissue Expression (GTEx) Consortium data to examine how a genetic variant potentially regulates its cis- and trans-target genes.
%, and the other is to profile the relationship between DNA methylation and gene expression in breast cancer patients, using copy number alteration as the instrumental variable.  

Network inference under the principle of Mendelian randomization (PMR) enables the study of causal relationships between two molecular phenotypes (e.g., gene expression) by leveraging genetic variants as instrumental variables and forming a trio. In practice, inference is often affected by numerous confounding variables, including other genes and demographic factors, making it essential to robustly account for confounders. However, existing methods have limited ability to accommodate high-dimensional confounding.
Here, we introduce Mendelian Randomization for Genomic Networks (MRGN), a regression-based method under the PMR for inferring causal relationships in genomic trios from individual-level data while accounting for many confounding variables. MRGN extends beyond the canonical mediation model to distinguish among multiple causal relationships. Through extensive simulations, we demonstrate improved inference accuracy, robustness, and computational efficiency over existing approaches. We further apply MRGN to data from the Genotype-Tissue Expression (GTEx) Consortium to investigate how genetic variants may regulate gene expression in cis and trans.

\setcounter{section}{0}
\section{Introduction}

An important goal in biology is to understand how genes regulate one another and how genetic variation influences gene expression \cite{camacho2018next, pingault2018using, xu2019machine, lecca2021machine}. Gene regulation involves complex networks of high dimensions, which may be represented by directed networks where directed edges represent causal (or regulatory) relationships. Although the ultimate goal in biology is to reconstruct gene regulatory networks at the genome level, it is often of interest to understand the causal relationship between a molecular phenotype (e.g. the expression level) of two genes, or between two molecular phenotypes (e.g., between gene expression and DNA methylation), which may be influenced by many other variables such as other genes, other biological processes, demographic variables, etc.

To learn the (directed) causal relationship between two phenotypes, we can apply the principle of Mendelian randomization (PMR) and leverage the information in genetic variants. The PMR assumes that genotypes in a natural population are randomly assigned, thus observational omics data can be viewed as being collected from a randomized experiment \cite{davey2014mendelian, aten2008using}. This implies that an edge connecting a genotype to a phenotype will always point to the phenotype, but not the other way around \cite{badsha2019learning, yazdani2016generating}. As a result, this constraint on edge direction makes it possible to break the symmetry of association and draw inferences about the regulatory structure among phenotypes \cite{davey2014mendelian, badsha2019learning, ludl2021comparison}.   

Methods based on the PMR framework differ fundamentally from classical Mendelian randomization (MR) methods. 
Traditional MR focuses on estimating causal effects and depends on three key assumptions 
\cite{davey2014mendelian, vanderweele2014mr, emdin2017mrreview}: (i) relevance, where the genetic variant is 
associated with the molecular phenotype; (ii) independence, where the variant is independent of unmeasured 
confounders due to random allele segregation; and (iii) exclusion restriction, where the variant influences downstream 
phenotypes only through the intermediate phenotype, and not through other pathways.
In contrast, our PMR-based framework does not require these assumptions to hold in full, nor is causal effect estimation 
its primary objective. Instead, following our earlier work \cite{badsha2019learning, badsha2021mrpc, kvamme2025cistrans}, we use the PMR 
as a structural constraint that limits edge directions in the network. Specifically, edges from molecular phenotypes to 
genetic variants are not allowed in the network. 

Many methods have been developed under assumptions similar to the PMR \cite{millstein2009cit, wang2014.network, yang2017identifying, wang2017efficient, badsha2019learning, wang2019high, howey2020bn, oh2024mr}, although some of these methods are interested mostly in the mediation model $V \rightarrow T_1 \rightarrow T_2$, where $V$ is the genetic variant, $T_1$ and $T_2$ are the variables of interest. In this model, the mediator $T_1$ relays the regulatory information from $V$ to the outcome $T_2$; for example, a genetic variant regulates the expression level of a transcription factor, and its reduced expression leads to its downstream target gene being under-expressed. However, other causal relationships may exist among the three variables. 

We previously developed MRPC, a method that goes beyond mediation and infers five mutually exclusive, basic causal models for a trio under the PMR (\autoref{fig:trio}a) \cite{badsha2019learning, badsha2021mrpc, kvamme2025cistrans}. They include a null model ($V \rightarrow T_1$ and no relationship between $T_1$ and $T_2$, denoted M0) and four non-null models: mediation (denoted M1), v-structure (or reverse causation; $V \rightarrow T_1 \leftarrow T_2$, denoted M2), conditional independence (or horizontal pleiotropy; $T_1 \leftarrow V \rightarrow T_2$, denoted M3), and a fully connected model ($T_1 \leftarrow V \rightarrow T_2$ and $T_1 \leftrightarrow T_2$, denoted M4). 
The true model behind M4 may be directed: the edge between $T_1$ and $T_2$ is $T_1 \rightarrow T_2$ or $T_1 \leftarrow T_2$, but these two directions result in the same likelihood for the trio and are therefore Markov equivalent. Without additional information, they are indistinguishable and only a bidirected/undirected edge can be inferred for M4.   


\begin{figure}[htbp]
	\centering
	\includegraphics[width = \textwidth]{./figures/Fig1_basic_models.pdf}
	\caption[Causal models for a genomic trio]{\textbf{Causal models for a genetic trio.}
	(a) The five basic causal models for a trio that may be inferred under the principle of Mendelian randomization (adapted from \cite{badsha2019learning}). 
	(b) A genomic trio includes a genetic variant ($V$), two molecular phenotypes ($T_1$ and $T_2$), and a set of confounders ($\mathbf{U}$). 
	%Regression coefficients from \autoref{eq:reg1} and \autoref{eq:reg2} correspond to the directed edges in the graph. 
	}
	\label{fig:trio}
\end{figure}

The PMR and the five basic causal models require that no confounding variables can explain away the directed edge between $T_1$ and $T_2$, since this edge is our main interest.  However, as mentioned earlier, many variables can influence the inference due to the complexity of gene regulation.  A confounding variable may influence the inference in different ways.  Yang et al.~\cite{yang2017identifying} proposed three types of confounding $C$ for the $T_1-T_2$ edge: as a confounder ($T_1 \leftarrow C \rightarrow T_2$), an intermediate variable ($T_1 \rightarrow C \rightarrow T_2$), or a common child ($T_1 \rightarrow C \leftarrow T_2$). Specifically, a confounder is a common parent to both $T_1$ and $T_2$ and can induce a false edge between $T_1$ and $T_2$ when unaccounted for.  This is similar to $T_1$ and $T_2$ being marginally dependent in M3.  
An intermediate variable acts like a mediator between $T_1$ and $T_2$ and can potentially explain away a true $T_1 \rightarrow T_2$ edge when this intermediate variable is accounted for in the inference, leading to a false negative conclusion.  By contrast, a common child can induce a false edge between $T_1$ and $T_2$ when accounted for, leading to a false positive conclusion.  This is because when the common child is accounted for in the inference, $T_1$ and $T_2$ become conditionally dependent, similar to their relationship in M2.  

Accounting for confounding variables is therefore crucial in causal network inference \cite{burgess2015network, yang2017identifying, shan2019identification, amar2021graphical, zhao2024adjusting}. Two basic approaches have emerged to address latent confounding between an exposure and outcome. In the first approach, candidate confounding variables are incorporated into the graph as additional nodes in the network and also treated in the same way as other nodes; our MRPC method and several other methods take this approach \cite{badsha2019learning, amar2021graphical}. However, this approach is computationally intensive, as the computational complexity for graph inference grows super exponentially with the number of nodes. The second approach focuses on identifying and accounting for confounders \cite{burgess2015network, yang2017identifying, shan2019identification}. This approach includes all the confounders as a whole in the tests for edge presence, and does not increase the size of the network.  For example, we can include the confounders when testing the $T_1 - T_2$ edge in a trio.  The trio remains a network of three nodes, but the tests are conditioned on these confounders.  However, these methods focus on the mediation model.

Herein we present Mendelian Randomization for Genomic Networks (MRGN), 
a novel machine learning method for inferring the causal model of a trio from individual-level data in the presence of potentially hundreds of confounding variables.  MRGN builds on MRPC and infers the five basic causal models.   We take the second approach to confounding described above and focus on accounting for confounders, and not other types of confounding variables.  This approach allows for the inclusion of many confounding variables all at once, while maintaining high power and fast computation times. We demonstrate the performance of MRGN and compare to other methods through extensive simulation studies.  We further assess this method in our study of the cis- and trans-regulatory modes of genetic variants, using the whole blood data from the Genotype-Tissue Expression (GTEx) Consortium \cite{gtex2020}. 

\section{Methods}
\subsection{The MRGN model}

The five causal models for a trio can each be described by a set of marginal and conditional dependence relations (\autoref{fig:trio}b; \autoref{tab:tests}), which may be encoded in the coefficients of two regression models that account for confounders, assuming that $T_1$ and $T_2$ are centered so that the intercept is zero:
\begin{align}
	T_1 &= \betaoneone V+\betaonetwo T_2+\boldGamma_1 \boldU+ \epsilon_1; \label{eq:reg1} \\
	T_2 &= \betatwoone V+\betatwotwo T_1+\boldGamma_2 \boldU+ \epsilon_2, \label{eq:reg2}
\end{align}
where $\beta$s are coefficients of interest, vector $\boldGamma_i$ the coefficients for the confounder vector $\boldU_i$, and $\epsilon_i$ the random error that follows a normal distribution with mean 0 and standard deviation $\sigma_i$, where $i=1, 2$.  

Our conditional independence tests test whether each $\beta$ is zero, as each $\beta$ represents the association given the presence of other variables.  For example, in the mediation model of $V \rightarrow T_1 \rightarrow T_2$, any two variables are dependent on each other, but $T_2$ and $V$ are conditionally independent given $T_1$.  These relations translate to (i) $\betatwoone$ being zero, which represents the association between $V$ and $T_2$ conditioned on $T_1$ and the confounders $\boldU$, and (ii) the other three $\beta$s being nonzero, which represent those associations not being explained away by other variables.  Similarly, models M0 and M3 are uniquely determined by conditional associations, for which we only need to conduct conditional independence tests.

In the v-structure model M2, consider the configuration $V \rightarrow T_1 \leftarrow T_2$. As the parents of $T_1$, nodes $V$ and $T_2$ are marginally independent but conditionally dependent given the presence of their child $T_1$.  Therefore, all the $\beta$s, which represent conditional associations, should be nonzero.  Similarly, all the $\beta$s are also nonzero in the fully connected model M4.  The distinction between M2 and M4 lies in the marginal associations: $V$ and $T_2$ are marginally independent in M2 but are marginally dependent in M4.  MRGN therefore performs additional tests on marginal dependence to distinguish these two models.
%For marginal association, we perform the standard test on Pearson correlation.

%the coefficient $\beta$ in the following univariate regression (\textcolor{red}{NEED}: which is implemented?):
%\begin{align}
%	T_i &= \alpha + \beta_{i1} V + \epsilon, \;\;\;\; \text{$i = 1$ or 2}. \label{eq:eq3}
%\end{align}

MRGN performs the standard $t$ tests on the regression coefficients for conditional dependence and the standard test on Pearson correlation for marginal dependence.  Since the number of tests for the edges of interest is at most six, we set the significance level for each test to be 0.01, which is equivalent to Bonferroni correction at a familywise type I error rate of 0.05.  When the data is not informative enough for any of the basic models to be inferred, MRGN categorizes the trio as ``other".  


\begin{table}[H]
	\centering
	\caption{\textbf{Conditional and marginal independence tests in MRGN.} Regression coefficients $\betaoneone$ and $\betaonetwo$ (from \autoref{eq:reg1}) and $\betatwoone$ and $\betatwotwo$ (from \autoref{eq:reg2}) are used to test conditional independence. Marginal independence can be tested through correlations $r_{V, T_1}$ and $r_{V, T_2}$.  The first three models each have two possible configurations.}
	\renewcommand{\arraystretch}{0.9}
	%\renewcommand{\arraystretch}{0.9} % Shrink space between rows
	\begin{tabular}{c|p{2.8cm}|cccc|cc} 
		\toprule[0.04cm]
		\multicolumn{2}{c|}{} & \multicolumn{4}{c}{Conditional Test} & \multicolumn{2}{|c}{Marginal Test}  \\ 
		\midrule
		
		\multicolumn{2}{c|}{Null Hypothesis} 
		& $T_1\indep V$ 
		& \small $T_1 \indep T_2$ 
		& \small $T_2\indep V$
		& \small $T_2\indep T_1$ 
		& \small $V\indep T_1$ 
		& \small $V \indep T_2$    \\
		
		\multicolumn{2}{c|}{Conditioned on} 
		& \small $T_2, {\bf U}$
		& \small $V, {\bf U}$
		& \small $T_1, {\bf U}$
		& \small $V, {\bf U}$
		& \small -- 
		& \small --    \\
		
		\multicolumn{2}{c|}{Parameter to Test} & $\betaoneone$ & $\betaonetwo$ & $\betatwoone$ & $\betatwotwo$ & $r_{V, T_1}$ &  $r_{V, T_2}$  \\
		\midrule
		
		Model & Configuration & \multicolumn{4}{c|}{} & \multicolumn{2}{c}{} \\ 
		\midrule[0.04cm]
		
		\multirow{4}{*}{M0} & \small $V \rightarrow T_1$; $T_2$ is singleton & $\neq 0$ & $= 0$ & $=0$ & $=0$ & -- &  --  \\ 
		
		& \small $V \rightarrow T_2$; $T_1$ is singleton &  $= 0$ & $= 0$ & $\neq 0$ & $= 0$ & -- & -- \\ 
		\midrule
		
		\multirow{2}{*}{M1} \small & $V \rightarrow T_1 \rightarrow T_2$ & $\neq 0$ & $\neq 0$ & $= 0$ & $\neq 0$ & -- & -- \\  
		
		& \small $V \rightarrow T_2 \rightarrow T_1$ & $ = 0$ & $\neq0$ & $\neq0$ & $\neq 0$ & -- & -- \\ 
		\midrule
		
		\multirow{2}{*}{M2} & $V\rightarrow T_1 \leftarrow T_2$ & $\neq0$ & $\neq 0$ & $\neq0$ & $\neq0$ & $\neq0$ & $=0$ \\ 
		
		& \small $V\rightarrow T_2 \leftarrow T_1$ & $\neq0$ & $\neq 0$ & $\neq0$ & $\neq0$ & $=0$ & $\neq 0$ \\  
		\midrule
		\multirow{1}{*}{M3} & \small $T_1 \leftarrow V \rightarrow T_2$ & $\neq 0$ & $= 0$ & $\neq0$ & $=0$ & -- & -- \\  
		\midrule
		\multirow{2}{*}{M4} & \small $T_1 \leftarrow V \rightarrow T_2$; $T_1 \longleftrightarrow T_2$ & $\neq0$ & $\neq 0$ & $\neq0$ & $\neq0$ & $\neq 0$ & $\neq 0$ \\  
		\bottomrule[0.04cm]
	\end{tabular}
	\label{tab:tests}
\end{table}


\subsection{Confounder selection}
We perform confounder selection before causal network inference, using a two-step procedure similar to that in the GMAC method~\cite{yang2017identifying}. First, we identify candidate confounders from a pool of potential confounding variables.  Confounders, which are defined as parents of $T_1$ and $T_2$, should be independent of $V$.  We therefore exclude variables that are marginally associated with $V$. 
%This filtering step is motivated by the assumption that true confounders are parents of the $T$ nodes and are therefore independent of each other and of $V$. In contrast, variables downstream of one or both $T$ nodes, such as common children or intermediates, may appear associated with $V$ due to collider bias or mediation effects, and should not be treated as confounders.
Next, we select confounders based on their statistical association with either $T_1$ or $T_2$. Similar to \cite{yang2017identifying}, we test whether the regression of both $T_1$ and $T_2$ on a candidate confounder is statistically significant.  Both steps allow for an additional step of the $q$-value method to control the false discovery rate (FDR) \cite{storey2003positive}.

In our R implementation, we substantially improve computational efficiency in two key ways. First, we implemented parallelization in both steps, allowing for faster evaluation when the pool of confounders is large. Second, we used the R package \textit{propagate} for correlation matrix computation by partitioning the problem into smaller, memory-efficient blocks \cite{spiess2018package}, which scales more efficiently to large genomic datasets.


\subsection{Simulation study}
%We examined the performance of MRGN in a wide range of settings and in comparison to similar methods; MRPC \cite{badsha2019learning} and GMAC \cite{yang2017identifying}. 
\subsubsection{Simulation setup}

Each simulated trio consists of a bi-allelic genetic variant $V$ with three genotypes and two genes $(T_1, T_2)$ connected under one of the five causal models. 
Nodes in a trio are generated in topological order. The genotype $V_i$ is simulated first under Hardy--Weinberg equilibrium:
\[
\Pr(V_i = 0) = (1-\theta)^2,\quad 
\Pr(V_i = 1) = 2\theta(1-\theta),\quad
\Pr(V_i = 2) = \theta^2,
\]
where the minor allele frequency $\theta$ is sampled uniformly from $(0.01, 0.5)$, such that roughly $10\%$ of trios contain a less common allele with a $<5\%$ frequency.
%For models with multiple configurations (M0--M2), we fix one representative direction:  
%\begin{itemize}
%    \item \textbf{M0:} $V \rightarrow T_1$; no edge between $T_1$ and $T_2$.
%    \item \textbf{M1:} $V \rightarrow T_1 \rightarrow T_2$.
%    \item \textbf{M2:} $V \rightarrow T_1 \leftarrow T_2$.
%\end{itemize}

Each remaining node $Y_j$ (a gene or a confounder) is sampled from
\[
Y_j \sim N\!\left(\sum_{i\in S_j} \beta_i Y_i,\;\sigma^2\right),
\]
where $S_j$ is the parent set of $Y_j$, which may include $V$, the other gene, or confounders.
The effect size of $V$ on its target genes is sampled uniformly from $[0.5, 1.5]$~\cite{yang2017identifying}. 
The effect size between $T_1$ and $T_2$ is sampled uniformly from $[0.5, 1]$. Residual variance is set to $\sigma = c \beta$,
where $\beta$ is the $T_1$--$T_2$ effect and $c$ sampled uniformly from $[0.3, 1.5]$.

For \textbf{M4}, the true edge between $T_1$ and $T_2$ can be in either direction. For simplicity, we fix all M4 trios to the configuration $T_1 \leftarrow V \rightarrow T_2$ with $T_1 \rightarrow T_2$.

For each trio, we further include three types of confounding variables: a varying number of confounders (between 1 and 50), one intermediate variable, and one common child. Effect sizes for these variables are sampled independently: confounder and intermediate effects from $[0.15, 0.5]$ and common-child effects from $[0.5, 1.5]$, each with a 50\% probability of being negative.

We simulate 1,800 trios by sampling from the five causal models with equal probability, and obtain approximately 360 trios per model. For each trio, we draw a distinct combination of effect sizes, noise level, and minor allele frequency. The sample size is fixed at $n = 670$, matching the whole-blood sample size of GTEx v8~\cite{gtex2020} in our real-data application.


\subsubsection{Performance metrics}

We evaluated three sets of recall (i.e., power) and precision to quantify the performance on the entire model of a trio, on individual edges and specifically on the $T_1-T_2$ edge.  %Then we calculated the mean (NEED to clarify) of each metric across all simulated trios.  

\begin{enumerate}
  \item On the causal model type:  a trio is considered correctly inferred when its inferred model type matches its true type, thus requiring all edges to be inferred correctly.
\begin{align}
	\text{Recall}_\text{model M} &= \frac{\text{Number of correctly inferred trios}}{\text{Total number of trios of the causal model M}}; \label{eq:rec_total}\\
	\nonumber \\
	\text{Precision}_\text{model M} &= \frac{\text{Number of correctly inferred trios}}{\text{Total number of trios inferred to be of the causal model M}}. \label{eq:prec_total}
\end{align}
  \item On individual edges: we first calculated recall and precision based on individual edges in a single trio.  Similar to our earlier work \cite{badsha2019learning}, we counted an edge with the wrong direction as a half correct edge.  
\begin{align}
	\text{Recall}_\text{edges} &= \frac{\text{Number of correct edges in inferred graph}}{\text{Total edges in true graph}}; \label{eq:rec_edges}\\
	\text{Precision}_\text{edges} &= \frac{\text{Number of correct edges In inferred graph}}{\text{Total edges in inferred graph}}.\label{eq:prec_edges}
\end{align}
We then calculated the mean metric across all trios.
  \item On the $T_1 - T_2$ edge (allowing for comparison with the GMAC method): the direct of this edge was ignored for this evaluation.  Across all trios,
\begin{align}
	\text{Recall}_\text{$T_1 - T_2$} &= \frac{\text{Number of trios with a correctly identified $T_1 - T_2$ edge}}{\text{Total trios with a true $T_1 - T_2$ edge}}; \label{eq:rec_tedge}\\
	\nonumber \\ 
	\text{Precision}_\text{$T_1 - T_2$} &= \frac{\text{Number of trios with a correctly identified $T_1 - T_2$ edge}}{\text{Total trios with an inferred $T_1 - T_2$ edge}}. \label{eq:prec_tedge}
\end{align}
\end{enumerate}


\section{Results}

\subsection{Simulation Results}
\subsubsection*{MRGN improves power and precision to detect edges in the presence of confounding}

Before applying MRGN, we examined two approaches to confounder selection: controlling the FDR at $5\%$ with the $q$-value method (denote ``CS-q"), and not correcting for multiple testing but instead setting the type I error rate of each individual test at $\alpha<0.01$ (denoted ``CS-$\alpha$").   We evaluated their ability to identify the true confounders, measured by precision and recall.  Here, recall is defined as the proportion of true confounders correctly identified and selected for the inference model, and precision is the proportion of true confounders selected out of all confounders selected.  Overall, both methods achieved recall greater than 90\% when the number of true confounders per trio was modest (e.g., fewer than 25; see \autoref{fig:SF4_Confounder_Selection_Performance}). However, CS-q showed reduced performance for trios with more true confounders, with average recall dropping to approximately 65\%. Conversely, as expected, CS-$\alpha$ achieved nearly perfect recall across all simulations, but at the expense of low precision.  Whereas the actual FDR (i.e., 1-precision) under CS-$\alpha$ was above 70\% and increased to above 90\% as the number of true confounders decreased, the actual FDR under CS-q was generally around 20\%, already much higher than the expected level, and increased to 50\% as the number of true confounders decreased.  This is likely due to the dependency among the tests: this step performs many tests involving the same variable or highly correlated variables.

\begin{figure}[htbp]
	\centering
	\includegraphics[width=\textwidth]{./figures/Figure_Confounder_Selection_Performance.pdf}
	\caption{\textbf{Precision and recall of confounder selection on simulated data by two multiple-testing correction methods: CS-q and CS-$\alpha$.} Each shape is a simulated data set.}
	\label{fig:SF4_Confounder_Selection_Performance}
\end{figure}

We applied MRGN with these two methods for confounder selection.  They led to qualitatively similar results in recall and precision on the inference of the trio (\autoref{tab:results}).  Both versions achieved substantially higher recall for the mediation, v-structure, and fully connected models than MRPC did. All these models contain a $T_1 - T_2$ edge, indicating the sensitivity of MRGN in detecting the $T_1 - T_2$ edge. On the other hand, compared to MRPC, this sensitivity led to reduced power to detect the null and conditional independence models, which have no $T_1 - T_2$ edge. Remarkably, MRGN also achieved equivalent or better precision for nearly all models, as well as similar performance between the two confounder selection methods, demonstrating robustness to high false positive rates in confounders.

\begin{table}[H]
	\caption{\textbf{Comparison of MRGN, MRPC and GMAC on simulated trios with confounders.} The oracle performance of MRGN using true confounders is shown to illustrate the upper bound due to noise in the data.  For the other columns, the highest precision and recall in each row is highlighted.  CS-q: confounder selection with the $q$-value method for FDR control at $q < 0.05$.  CS-$\alpha$: confounder selection with the type I error rate set at $\alpha < 0.01$ and no FDR control.  MRPC cannot handle too many variables in reasonable runtime; we therefore did not run MRPC with the many confounders identified from CS-$\alpha$. Rows 1-5 show model type-based performance, row 6 for all three edges in a trio, and the last row  for only the $T_1-T_2$ edge.  $\ast$ indicates that similar values are produced by the GMAC method.}
%, and the last row gives the performance across all edges in the network using the weighting scheme described in the Section ``Performance metrics".}
	\renewcommand{\arraystretch}{0.9} % Shrink space between rows
	\begin{tabular}{l|ll|ll|ll|ll}
		\toprule
		 & \multicolumn{2}{p{2.8cm}|}{\centering \bf MRGN} &  \multicolumn{2}{p{2.8cm}|}{\centering \bf MRGN} & 
		\multicolumn{2}{p{2.8cm}|}{\centering \bf MRGN} & 
		\multicolumn{2}{p{2.8cm}}{\centering \bf MRPC} \\
		 & \multicolumn{2}{p{2.8cm}|}{\centering \bf (True Confs.)} &  \multicolumn{2}{p{2.8cm}|}{\centering \bf (CS-q)} & 
		\multicolumn{2}{p{2.8cm}|}{\centering \bf (CS-$\alpha$)} & 
		\multicolumn{2}{p{2.8cm}}{\centering \bf (CS-q)} \\		
		& Recall & Prec. & Recall & Prec. & Recall & Prec. & Recall & Prec. \\ 
		\midrule[0.04cm]
		Null Model ({\bf M0}) & 0.562  & 0.995  & 0.452  & \bf 1.000  & 0.477  & 0.951  & \bf 0.836  & 0.532  \\
		Mediation ({\bf M1}) & 0.982  & 0.516  & 0.877 & \bf 0.438  & \bf 0.883  & 0.429  & 0.097  & 0.147  \\
		V-structure ({\bf M2}) & 0.919  & 0.985  & \bf 0.892  & 0.925  & 0.842  & \bf 0.968  & 0.377  & 0.443  \\ 
		Cond. Indep. ({\bf M3}) & 0.515  & 0.989  & 0.283  & \bf 0.910  & 0.305  & 0.872  & \bf 0.609  & 0.828  \\ 
		Fully Connect. ({\bf M4}) & 0.600  & 0.566  & \bf 0.501  & 0.425  & 0.464  & 0.419  & 0.217  & \bf 0.592  \\ 
		All Edges  & 0.945  & 0.893  & \bf 0.931  & 0.864  & 0.911  & 0.857  & 0.770  & \bf 0.892  \\
		$T_1-T_2$ & 1.000  & 0.767  & \bf 0.994$^*$  & 0.703$^*$  & 0.9787 & 0.7071 & 0.7884 & \bf 0.916 \\
		\bottomrule
		%\multicolumn{9}{c}{\footnotesize $\ast$ same values observed for GMAC}
	\end{tabular}
	\label{tab:results}
\end{table}

At the edge level, MRGN recovered more true edges than MRPC did, while maintaining a high degree of precision. 
MRGN and GMAC performed similarly well in recall and maintained consistently high recall for the $T_1 - T_2$ edge, even under substantial confounding, with precision remaining relatively unchanged (\autoref{fig:main_t1_t2_inf}). This is because the GMAC test for mediation is a special case of our method. The similar performance on the $T_1 - T_2$ edge further shows that the notion of ``mediation" detected by GMAC is broad and an inferred mediation edge by GMAC may correspond to different underlying causal models.  Additionally, recall and precision for all the edges in a trio remained high, indicating that MRGN can accurately infer most edges, except in cases with many ($>35$) confounders, where recall declined but still exceeded that of MRPC (\autoref{fig:main_all_edge}).

\subsubsection*{The impact of confounder selection on causal network inference}

When we examined how confounder selection affected causal inference performance, we found no clear relationship between the number of falsely included confounders and MRGN’s ability to accurately infer causal edges (\autoref{fig:SF7_Confounder_Selection_Impact_On_Inference}a). In contrast, falsely omitted confounders had a strong negative impact on causal inference performance (\autoref{fig:SF7_Confounder_Selection_Impact_On_Inference}b). 
%Overall, we observed little difference in causal inference performance between the two confounder selection approaches in our simulation study. This may be because only a small number of trios were simulated with a large number of confounding variables (e.g., more than 35), limiting our ability to detect any differences between the two methods. Although CS-$\alpha$ is more liberal, it may be preferable in scenarios where many potential covariates are expected for a trio.

\begin{figure}[H]
	\centering
	\includegraphics[width = \textwidth]{./figures/SF7_Confounder_Selection_Impact_On_Inference.pdf}
	\caption{\textbf{Relationship between confounder selection and causal inference performance measured as precision and recall of edge prediction.} 
	(a) Edge-based performance of MRGN on simulated trios as a function of the number of falsely included confounders.
	(b) Edge-based performance of MRGN on simulated trios as a function of the number of true confounders falsely omitted from the analysis.}
	\label{fig:SF7_Confounder_Selection_Impact_On_Inference}
\end{figure}

We further observed substantially higher actual type I error rates for the $T_1 - T_2$ edge under the GMAC method 
than those reported by \cite{yang2017identifying} (\autoref{tab:T1.T2.Edge.Results.Combined.csv}). 
Their actual type I error rates were consistent with the significance level, such as $\alpha=0.01$ or 0.05. 
However, the actual type I error rates in our simulation using GMAC are 0.73 at $\alpha=0.05$ and 0.63 at $\alpha=0.01$.
Using MRGN, our actual type I error rates are also around 0.6 (with $\alpha=0.01$ in causal network inference).
%(NEED to check which result? supplementary_tables/T1.T2.edge.results.csv). 
We attribute these differences to several distinctions between our simulation framework and theirs. 

First, the number of confounding variables is much larger in our simulation than in theirs.  
We simulated a substantially larger number of confounders (between 1 and 50) for each trio. 
We further aggregated the confounding variables from all the trios for confounder selection, 
resulting in a large candidate pool of more than $10,000$ variables. 
In contrast, \cite{yang2017identifying} simulated a much smaller pool of $350$ confounders, 
yielding an average of approximately two confounders per trio, with many trios having no confounders. 
Second, we simulated two additional confounding variables for each trio: one common child and one intermediate variable.  
These variables were also included in the candidate covariate pool. 
The simulations in \cite{yang2017identifying}, however, treated each type of confounding variable separately.

To validate the error rates reported by \cite{yang2017identifying}, we used their simulation design to examine 
the effect of confounders on type~I error. We simulated 100 additional trios for each of seven scenarios containing 
2, 4, 6, 8, 10, 12, or 15 confounders, excluding common child and intermediate variables.  We aggregated all the confounders in the 
candidate pool, and applied GMAC with CS-q for confounder selection. The actual type I error rate under two true confounders is comparable to  
that reported by \cite{yang2017identifying} (\autoref{fig:SF9_GMAC_Validation_Summary}), whereas all other scenarios showed that the actual 
type I error rate scaled with the number of confounders and the size of the candidate pool. These results demonstrate that more confounding 
leads to markedly inflated false positive $T_1 - T_2$ edges and decreased power to detect the null (M1) and conditional independence (M3) models by both GMAC and MRGN.

\subsubsection{Runtime comparison}

Since our analysis involves two major steps: confounder selection and causal network inference, we report the computation time for each step separately. We conducted a comparative analysis of the time required for confounder selection across varying sizes of the confounder covariate pool (100, 500, or 1000 covariates) and different numbers of trios for which confounders needed to be identified (500 or 1000), resulting in a total of six replicates. Notably, in the most intensive case 1000 covariates and 1000 trios, our method accomplished confounder selection within 16 minutes (\autoref{fig:SF1_comp_times_for_conf_selection_updated}). This represents a notable gain in computation efficiency compared to the implementation provided by GMAC, which took approximately 28 minutes. Although our confounder selection procedure is similar to that of \cite{yang2017identifying}, our implementation is more efficient due to the use of the R package propagate \cite{spiess2018package}, which enables fast correlation calculations between genes in each trio and the candidate covariate pool, significantly reducing the time required for confounder selection.

\begin{figure}[H]
	\centering
	\includegraphics[width=4.4in]{./supplementary_figures/SF1_comp_times_for_conf_selection_updated.pdf}
	\caption{\textbf{Comparison of computation times between our implementation of the CS-q method for confounder selection and that of the GMAC method \cite{yang2017identifying}}. Computation times are shown across different numbers of input trios and sizes of the candidate confounder pool.}
	\label{fig:SF1_comp_times_for_conf_selection_updated}
\end{figure}

For the step of causal network inference, we evaluated the runtime on all the 1,800 trios.  MRGN is substantially faster than both MRPC and GMAC (\autoref{tab:time_comp}). The median time to analyze a trio using the same set of confounders was 0.0065 seconds for MRGN, compared to 0.3557 seconds for MRPC and 0.7088 seconds for GMAC. This makes MRGN over 50 times faster than MRPC and more than 100 times faster than GMAC, on average.

The extended inference time for MRPC is largely due to the ADDIS method used for FDR control \cite{kvamme2025cistrans}, which was particularly slow for some trios. When MRPC was applied using the more conservative LOND method for FDR control \cite{badsha2019learning}, its runtime became comparable to that of MRGN. For GMAC, we note that our comparison used the default single-threaded implementation provided in its R package. Enabling the parallelization option available in GMAC can substantially reduce its computational time.

\begin{table}[H]
	\caption{\textbf{Inference times for MRGN, MRPC, and GMAC on simulated trios}. We used the ADDIS method for FDR control in MRPC \cite{kvamme2025cistrans}. MRGN shows a substantially faster median inference time per trio compared to both MRPC-ADDIS and GMAC, demonstrating its computational efficiency in the presence of confounding.}
	\label{tab:time_comp}
	\centering
	\begin{tabular}{lcc}
		\toprule
		\textbf{Method} & \textbf{Median Time (s)} & \textbf{IQR (s)} \\
		\midrule
		MRGN & 0.011 & 0.005 \\
		MRPC-ADDIS & 0.460 & 9.012 \\
		GMAC & 1.537 & 1.483 \\
		\bottomrule
	\end{tabular}
\end{table}

%\subsubsection{More confounding contributes to an inflated type I error rate}




\subsection{Application to GTEx data}

The GTEx consortium has mapped a large number of genetic variants as expression quantitative trait loci (eQTLs), which are associated the expression for most genes in the human genome across many human tissue types \cite{gtex2020}. Recent  studies have focused on identifying cis-target genes that mediate the relationship between the eQTL and its associated trans-target genes \cite{yang2017identifying}.  We previously have applied the MRPC method to all the GTEx tissues to detect diverse causal models for such trios: an eQTL and its associated cis- and trans-target genes, where the genes are protein coding genes or long non-coding RNAs (lncRNAs) \cite{kvamme2025cistrans}.  

We implement the preprocessing and normalization pipeline from Section 2.2 to construct trios for WholeBlood tissue from version 8 of GTEx. Each trio consists of an eQTL and its highest association ``cis" and ``trans" target genes.  We define cis genes as genes less than one megabase from the eQTL and trans genes as genes residing more than one megabase from the eQTL or on a different chromosome. As noted in Section 2.2 individual trios may share the same eQTL or the same gene, and the gene may be a cis-gene in one trio but a trans gene in another trio. 

In GTEx version 8, the number of available samples for whole blood is 670. After preprocessing, 8,823 trios were identified. To eliminate bias induced by mismapping transcripts, we restricted our final set of trios to protein coding genes or long non-coding RNA genes \cite{yang2017identifying}. This additional filtering step reduces the number of trios to 3,248.  Many of eQTLs provided by the GTEx consortium contain missing values. Since the GMAC algorithm is unable to handle missing values, we used multiple correspondence analysis implemented in the \textit{missMDA} package in $R$ \cite{josse2016missmda} to impute the missing genotypes of individuals using the collection of unique eQTLs from all GTEx trios analyzed for whole blood. 

\subsubsection{Confounding variable selection and filtering for GTEx trios}

We performed principal component analysis (PCA) on the PEER-normalized, genome-wide gene expression data to derive $670$ 
principal components (PCs) as candidate confounders.  
Each PC summarizes the regulatory effects of a large number of genes across the genome. 
Because each trio represents a small component of a larger regulatory network, PCs that correlate with a trio may capture 
potential confounding effects arising from pleiotropy. 
%We performed confounder selection as described in Confounding Variable Selection and applied common-child and intermediate-variable filtering to obtain a final set of confounders for each trio. 
%Following our findings from simulation, 
We applied MRGN using CS-$\alpha$ at a more relaxed criterion of $\alpha < 0.05$ than that used 
in the simulation, in order to include more potential confounders. 
%While the simulation study 
%used a more stringent cutoff of $\alpha < 0.01$, we relaxed this threshold here because the more stringent cutoff produced 
%nearly the same set of confounders as CS-q. 
For each trio, we identified a median of $51$ confounders (SD = $6.38$) (\autoref{fig:main_GTEx}a). 
We also included sex, platform, and PCR batch as additional, known confounders for each trio.


\subsubsection{MRGN identified more mediation and conditional independence trios}
%We applied MRGN combined with CS$-\alpha$ to $3{,}248$ eQTL trios for whole blood in GTEx. We compared the results from MRGN with those from MRPC (using the ADDIS method to control the overall FDR of conditional tests) and GMAC \cite{badsha2019learning,yang2017gmac}. Principal component analysis (PCA) was performed across all genes in whole blood to derive $670$ principal components, which served as the pool of candidate confounding variables. For each trio, we identified a median of $51$ confounders (SD = $6.38$) (\autoref{fig:main_GTEx}). 

\begin{figure}[H]
	\centering
	\includegraphics[width = \textwidth]{./figures/MF3_GTEx.model.and.t1t2.edge.bargraphs.pdf}
	\caption[Performance of MRGN on GTEx human whole blood trios]{
		\textbf{Comparative analysis of MRGN, MRPC \cite{Badsha2019} and GMAC \cite{yang2017identifying} on GTEx human whole blood trios.} 
		We applied MRGN, GMAC, and MRPC (with ADDIS for FDR control) to genomic trios. 
		(a) Histogram of the number of Principal Components (PCs) selected as confounders per trio.
		(b) Count of inferred causal models (M1--M4 and Other) comparing MRGN and MRPC. 
		(c) Counts of trios with the $T_1 - T_2$ edge inferred as absent or present by the three methods. 
		(d) Counts of two types of mediation models (via cis-gene vs. trans-gene) inferred by MRGN and GMAC. 
	}
	\label{fig:main_GTEx}
\end{figure}

MRGN identified $911$ trios as the null model, $91$ as mediation, $10$ as v-structures, $2,062$ as conditionally independent, 
104 as fully connected, and $70$ trios as ``other" (\autoref{fig:main_GTEx}b), although the overall number of inferred $T_1-T_2$ edges 
was similar to that from the other methods (\autoref{fig:main_GTEx}c). Comparing with our previous analysis using MRPC \cite{kvamme2025cistrans}, 
we found that MRGN identified more than twice as many cases of mediation and $173$ additional cases of conditional independence. 
Similar to MRPC, MRGN also distinguishes two types of mediation (\autoref{fig:main_GTEx}d): (1) cis-gene mediation $V \rightarrow T_1 \rightarrow T_2$, 
where the cis-gene mediates the signal from genetic variant to the trans-gene, and (2) trans-gene mediation $V \rightarrow T_2 \rightarrow T_1$, 
where the trans-gene acts as the mediator between the genetic variant and the cis-gene. Among the $91$ trios identified as mediation by MRGN, 
$69$ (75.5\%) had the cis-gene as the mediator and $22$ (24.5\%) had the trans-gene as the mediator (\autoref{tab:S7_GTEx_comp_MRGN_MRPC_GMAC}). 
Compared to MRPC, MRGN identified not only more mediation trios overall, but also a higher proportion of cis-gene mediation. 

However, only $8$ cis-gene mediation trios and $6$ trans-gene mediation trios were identified by both methods.
We investigated trios where MRGN inferred mediation but MRPC did not and found that these discrepancies were primarily driven by differences in the sets of selected confounders (see Examples 1 and 2 in \autoref{txt:supptext_document}). In two illustrative cases, MRGN selected 40 and 43 confounders, whereas MRPC selected only three and none. When we reran MRGN using the same confounder set and inference model as MRPC, both methods produced identical results, confirming that accurate identification of confounders is essential for consistent causal inference.


\subsubsection{MRGN is robust to model misspecification for inference of mediation}

Based on our simulation and application findings, we conducted a secondary analysis to investigate how model misspecification affects the inference of the mediation model, since this causal relationship has been the focus of many studies. In particular, we examined how inaccuracies in the inference model, particularly in terms of confounding variables, influence the power to detect mediation on the $T_1 - T_2$ edge. We base our simulations on $192$ Whole Blood trios that were positive for mediation under the GMAC mediation test. We consider four simulation scenarios: (i) the inference model included true confounders and a few false positive confounders; (ii) the inference model included most true confounders, but missed a few true confounders with weak effects; (iii) the inference model included only the true confounders; and (iv) the inference model included only a few true confounders with strong effects, and missed many true confounders with moderate to strong effects (See \autoref{txt:supptext_document} for details and methods).  This investigation is similar to the simulation study that examined the impact of confounder selection on model inference, but further takes into consideration the effect sizes of confounders here.

Recall is around $93\%$ in each of the first three scenarios, indicating that when the inference model contains most of the influential confounding variables, the correct inference for the $T_1 - T_2$ edge can typically be recovered. Indeed, the first two scenarios represent relatively minor changes in the confounding variables used in the inference model, either by omitting or including weak confounders. On the other hand, even the correct model failed on a small subset of trios due to noise in the data.
In the final scenario, however, the inference model omitted an average of $28$ confounding variables (SD $\approx 7$), many of which had moderate to strong confounding effects. Recall is reduced to as low as $70\%$, highlighting the critical need to retain high-impact confounders for accurate causal network inference.

\section{Discussion}
In summary, we have developed MRGN to improve the power and accuracy of detecting key causal relationships in diverse genomic trios, particularly in the presence of multiple confounding variables. Our method adopts an extended interpretation of the PMR \cite{badsha2019learning} and leverages conditional and marginal hypothesis testing to identify five basic causal trio structures. A major advantage of our approach is the reduction of the multiple testing burden; by adjusting for selected confounders within a regression framework, MRGN can incorporate a large number of confounders without the need to infer a graph larger than three nodes.  

Our simulation results further underscore the critical importance of confounder selection in causal inference. Omitting key confounders drastically reduces the power to detect mediation, leading to Type II errors and the loss of regulatory insights. Conversely, the inclusion of false-positive confounders has a negligible impact, particularly when their effects are weak. This suggests that in practical applications, a liberal strategy favoring the inclusion of many potential confounders is preferable. Ultimately, MRGN balances this requirement without excessive computational cost, making it a robust tool for causal network inference.  However, when the sample size is not much larger than the number of identified confounders, including many false confounders can substantially increase the number of regression coefficients, leading to unreliable estimation.  We are currently developing extensions of MRGN to address issues with a small sample size.

Constraint-based methods such as MRPC initialize learning by filtering edges from a fully connected graph structure. After removing edges with weak marginal relationships, the remaining edges are tested sequentially by conditioning on an increasing neighborhood of nodes \cite{badsha2019learning, spirtes2000causation}. 
%Because the number of tests required to determine the existence of an edge is unknown a priori, MRPC utilizes sequential FDR control \cite{tian2019addis}. Consequently, the testing burden is heavily dependent on graph connectivity, signal strength, and the total number of nodes.
Critically, MRPC treats confounding variables as additional nodes within the graph. When confounders are included, the number of conditional tests increases significantly, and only edges with strong supporting evidence survive. Thus, MRPC inherently infers a larger network than MRGN for the same biological question. Additionally, the adaptive FDR can result in significance thresholds that deviate substantially from the fixed $\alpha = 0.01$ employed by MRGN. As demonstrated by our simulations, this causes MRPC to act more conservatively in the presence of confounding variables.

While our simulations and real data analysis demonstrate the flexibility of MRGN in robustly estimating causal relationships, the method is currently limited to trios containing a genetic variant that acts as an instrumental variable. Similar to MRPC, MRGN relaxes strict instrumental variable assumptions by empirically validating the association between the variant and the phenotype(s). We are currently extending this approach to develop a general MRGN algorithm capable of accommodating larger networks and addressing other classes of confounding variables.

\end{doublespace}

\section*{Data access}
Data from The Genotype-Tissue Expression (GTEx) Project version 8 were obtained from dbGaP (\url{http://www.ncbi.nlm.nih.gov/gap}) through 
accession number phs000424.v8.pht002741.v8.p2.  Open-source R packages used in this manuscript include MRGN (version 1.3.7.3);  \url{https://github.com/audreyfulab/MRGN}, MRPC (version 3.1.0); \url{https://cran.r-project.org/web/packages/MRPC/index.html}, and GMAC (version 3.1; 
\url{https://cran.r-project.org/web/packages/GMAC/index.html}.  R code for simulations and applications is available at \url{https://github.com/audreyfulab/MRGN_Manuscript_Repo}.  Detailed results generated by our analyses are included in the supplementary 
material.

\section*{Competing interest statement}
The authors declare no competing interests.


\section*{Acknowledgments}
We acknowledge the extensive computing support provided by the Research Computing and Data Services from the Institute for Interdisciplinary Data Sciences at the University of Idaho.

This research is supported by the National Institutes of Health (NIGMS P20GM104420).

The Genotype-Tissue Expression (GTEx) Project was supported by the Common Fund of the Office of the Director of the National Institutes of Health. Additional funds were provided by the NCI, NHGRI, NHLBI, NIDA, NIMH, and NINDS. Donors were enrolled at Biospecimen Source Sites funded by NCI/Leidos Biomedical Research, Inc.\ subcontracts to the National Disease Research Interchange (10XS170), Roswell Park Cancer Institute (10XS171), and Science Care, Inc. (X10S172). The Laboratory, Data Analysis, and Coordinating Center (LDACC) was funded through a contract (HHSN268201000029C) to the The Broad Institute, Inc. Biorepository operations were funded through a Leidos Biomedical Research, Inc.\ subcontract to Van Andel Research Institute (10ST1035). Additional data repository and project management were provided by Leidos Biomedical Research, Inc.(HHSN261200800001E). The Brain Bank was supported supplements to University of Miami grant DA006227. Statistical Methods development grants were made to the University of Geneva (MH090941 \& MH101814), the University of Chicago (MH090951,MH090937, MH101825, \& MH101820), the University of North Carolina - Chapel Hill (MH090936), North Carolina State University (MH101819), Harvard University (MH090948), Stanford University (MH101782), Washington University (MH101810), and to the University of Pennsylvania (MH101822). The datasets used for the analyses described in this manuscript were obtained from dbGaP at \url{http://www.ncbi.nlm.nih.gov/gap} through dbGaP accession number phs000424.v8.pht002741.v8.p2.

\section*{Author Contributions}
J.K and A.Q.F designed the study, J.K developed the tools and analyzed the data, and J.K and A.Q.F wrote the manuscript.


\bibliography{./references_sorted.bib} 
\bibliographystyle{unsrt}

\newpage
\section*{Tables}





\section*{Supplementary Material}

\newpage
\subsection*{Supplementary figures}
%\begin{suppfigure}[H]
%	\centering
%	\includegraphics[width = \textwidth]{./supplementary_figures/Confounding Vars2.pdf}
%	\caption{\textbf{The three types of confounding variables to consider in causal inference of genetic trios}. A confounder (left) is a gene which regulates both molecular phenotypes, an intermediate variable (center) is a gene which mediates part of the effect of the genetic variant and its target, and a common child (right) is a shared target of both molecular phenotypes.}
%	\label{fig:confounding}
%\end{suppfigure}


\begin{suppfigure}[H]
	\includegraphics[width = \textwidth]{./supplementary_figures/MF1_MRGN.GMAC.MRPC.all.params.t1t2.pdf}
	\caption{\textbf{Precision-recall curves for predicting the $T_1 - T_2$ edge across multiple simulation parameters}. (A) the number of simulated hidden confounders (B) Minor allele frequency of the genetic variant, (C) The signal strength of the $T_1 - T_2$ edge, and (D) The standard deviation of the regression residuals.}
	\label{fig:main_t1_t2_inf}
\end{suppfigure}

\begin{suppfigure}[H]
	\includegraphics[width = \textwidth]{./supplementary_figures/SF_MRGN.MRPC.all.params.edge.based.pdf}
	\caption{\textbf{Precision and recall curves of edge-based prediction performance for MRGN, and MRPC with confounder selection}. (A) the number of simulated hidden confounders (B) Minor allele frequency of the genetic variant, (C) The signal strength of the $T_1 - T_2$ edge, and (D) The standard deviation of the regression residuals.}
	\label{fig:main_all_edge}
\end{suppfigure}


\begin{suppfigure}[H]
	\centering
	\includegraphics[width = \textwidth]{./supplementary_figures/SF8_Trio_Skew_Stat_Distribtuions.pdf}
	\caption{\textbf{Skewness and kurtosis for cis and trans genes by allele frequency}. Stacked histograms showing the sampling distributions of sample skewness and excess kurtosis for cis and trans genes in GTEx trios. Results are segregated by the minor allele frequency of the genetic variant, highlighting differences between common and rare alleles.}

	\label{fig:SF8_Trio_Skew_Stat_Distribtuions}
\end{suppfigure}

%\begin{suppfigure}[H]
%	\centering
%	\includegraphics[width = \textwidth]{./supplementary_figures/trio_number_730.pdf}
%	\caption{\textbf{Example of a trio with a rare variant acting as an outlier.} In this trio, a single individual homozygous for the minor allele drives non-normality in trans-gene expression relative to the population. This individual (yellow point) also shows elevated expression of the cis-gene. Initially, MRGN inferred model M4 for this trio, while GMAC inferred no mediation. MRGN with permutation recovers the same inference (the null model) as GMAC.}
%	\label{fig:trio_number_730}
%\end{suppfigure}



\begin{suppfigure}[H]
	\centering
	\includegraphics[scale =0.5]{./supplementary_figures/SF9_GMAC_Valid_TypeI_error.pdf}
	\caption[Comparison of GMAC and MRGN type I error rates for $T_1-T_2$ detection under confounding]{Results of simulations investigating type I error rate (false mediation) under the null model. Trios were simulated under the null model for each of seven scenarios: 2, 4, 6, 8, 10, 12 or 15 confounders, excluding common child and intermediate variables. The GMAC method was applied to trios in each scenario using CS-q confounder selection.}
	\label{fig:SF9_GMAC_Validation_Summary}
\end{suppfigure}


\subsection*{Supplementary tables}


\begin{supptable}[H]
	\caption[Simulated trios and inference results for MRGN, MRPC-ADDIS, and GMAC]{
		A complete list of simulated trios, their simulation parameters, and
		inference results under MRGN, MRPC-ADDIS, and GMAC.
	}
	\label{tab:ST_all_results_simulations.csv}
\end{supptable}

\begin{supptable}[H]
	\caption[Inference for the $T_1-T_2$ edge in simulated trios byr MRGN, MRPC-ADDIS, and GMAC]{
		Different confounder selections were applied in MRGN and GMAC.  The table summarizes the number of trios inferred 
		to detect (or miss) the presence (or absence) the $T_1-T_2$ edge, as well as recall, precision, and actual type I and II error rates.
	}
	\label{tab:T1.T2.Edge.Results.Combined.csv}
\end{supptable}

\begin{supptable}[H]
	\caption[GTEx trios and inference results for MRGN, MRPC-ADDIS, and GMAC]{
		All analyzed GTEx whole-blood trios, including genetic annotations
		(locus, variant ID, gene name, etc.) and inference results from
		MRGN, MRPC-ADDIS, and GMAC.
	}
	\label{tab:ST_GTEx_all_trios_master.csv}
\end{supptable}


\begin{supptable}[H]
	\caption[Comparison of MRGN, MRPC, and GMAC on GTEx trios]{This table compares MRGN, MRPC, and GMAC on results from application to GTEx whole blood tissue. Results compared between MRGN with CS-$\alpha$ confounder selection to MRPC and GMAC. }
	\label{tab:S7_GTEx_comp_MRGN_MRPC_GMAC}
\end{supptable}




\subsection*{Supplementary text}

\begin{supptext}[H]
	\caption{Detailed supplementary text to explain the four simulation scenarios presented under the section ``MRGN is robust to model misspecification in causal mediation". This text also includes supplementary examples with output illustrating differences in causal inference between MRGN, GMAC, and MRPC on selected GTEx and simulated trios.}
	\label{txt:supptext_document}
\end{supptext}

\end{document}



























